% ------------------------------------------------------------------
% ============================================================
%  2.4  The Mapping Problem
% ============================================================
\section{The Mapping Problem}
\label{sec:bg-mapping}

Given a DNN workload expressed as a loop nest
(Section~\ref{sec:bg-workloads}) and a spatial architecture described as
a hierarchy of levels (Section~\ref{sec:bg-hw-accel}), a
\textbf{mapping} is a complete specification of how the computation is
partitioned, scheduled, and assigned to the hardware.
Concretely, a mapping answers three questions:
\begin{enumerate}
  \item \textbf{Where} does each operand tile reside in the memory
    hierarchy?
  \item \textbf{When} and \textbf{how} does data move between levels,
    and which multicast or reduction opportunities are exploited?
  \item \textbf{Which PE} computes \textbf{which MAC}, and in what
    temporal order?
\end{enumerate}

The set of all mappings that satisfy the resource and user constraints
of a given workload--architecture pair is called the
\textbf{map-space}~$\mathcal{MS}$.
The \textbf{mapping problem} is to find a mapping
$m^{*}\!\in\mathcal{MS}$ that minimizes a target metric, typically
energy, latency, or the energy-delay product (EDP):
\begin{equation}
\label{eq:mapping-problem}
m^{*} = \arg\min_{m\in\mathcal{MS}} \; \text{EDP}(m).
\end{equation}

\noindent
This optimization is difficult for two reasons.
First, the map-space is combinatorially large
(Section~\ref{sec:bg-mapspace} quantifies this).
Second, energy and latency are dominated by \emph{data movement}
rather than by computation: accessing off-chip DRAM costs hundreds of
times the energy of a PE-local register read, so the objective function
is highly sensitive to how effectively a mapping exploits data reuse.


% ------------------------------------------------------------
\subsection{Three Degrees of Freedom}
\label{sec:mapping-dof}
% ------------------------------------------------------------

Every mapping is fully characterized by three orthogonal choices made
\emph{independently at each level} of the architecture hierarchy.

% ---- TILING ----
\subsubsection{Tiling}
\label{sec:tiling}

To fit a workload's tensors into the limited on-chip memories of a
spatial architecture, the loop nest is \emph{tiled}: each loop
dimension~$d$ with extent~$|d|$ is split into sub-ranges distributed
across the levels of the hierarchy.
Formally, the extent is factorized into its prime factors, and each
prime-power copy is assigned to exactly one level.
The product of the factors assigned to a level and all levels below it
determines the \emph{tile size} at that level, the number of elements
along dimension~$d$ that the level must store simultaneously for an iteration.

Tiling choices directly determine the storage requirements at each
memory level.
A mapping is \emph{valid} only if every tile fits within the capacity
of the memory level that must hold it.

\paragraph{Example.}
Consider a dimension $M = 128 = 2^7$ in a three-level hierarchy
(DRAM -- GlobalBuffer -- Register).
One valid tiling assigns factors~$2^3$ to DRAM, $2^2$ to the
GlobalBuffer, and $2^2$ to the Register, giving tile sizes of
$128$, $16$, and~$4$ respectively.
Each register tile of 4~elements is reused across the 4 iterations
at that level before a new tile is loaded from the GlobalBuffer; each
GlobalBuffer tile of 16~elements is reused across 8 iterations before
fetching from DRAM.

% ---- PARALLELISM ----
\subsubsection{Parallelism Strategy}
\label{sec:parallelism}

Identically to tiling: the prime-factor copies of a
dimension are distributed across all levels, including the spatial
ones.
The product of factors at a spatial level gives the number of
instances executing in parallel along that dimension.

The parallelism strategy interacts with data reuse through
\emph{multicast} and \emph{reduction}.
When a fanout level unrolls a dimension orthogonal to an operand,
all instances need the same data tile.
If the hardware supports spatial multicast, a single read from
the outer memory feeds all instances, multiplying the effective
reuse by the fanout mesh size.
Conversely, when the unrolled dimension is coupled to the output,
multiple instances produce partial sums for the same output tile;
spatial reduction support allows these to be accumulated in a single
write.

When the unrolled dimension is \emph{affine} for an operand,  adjacent instances
access \emph{overlapping} but not identical input tiles.
With \emph{selective multicast} support, the architecture can read
the union of all required tiles once and distribute the appropriate
sub-tiles to each instance.
Without it, each instance issues its own independent read, multiplying
bandwidth demand.

% ---- LOOP ORDERING ----
\subsubsection{Loop Ordering (Dataflow)}
\label{sec:loop-ordering}

At each memory level, the loops assigned to that level can be
\emph{permuted} without affecting correctness.
The chosen permutation constitutes the level's \textbf{dataflow} and
is the primary lever that controls which reuse mechanisms are
activated.

On a memory level, the operands that remain \emph{constant} across
the innermost iterations are kept stationary in that memory, avoiding
redundant reads.
The innermost loops therefore determine which operand is stationary.

The following reuse patterns arise from the interaction between loop
ordering and the coupling structure:

\begin{enumerate}
  \item \textbf{Stationarity (temporal reuse).}
    When the innermost iterated dimension at a memory level is
    \emph{orthogonal} to an operand
    (cf.~Definition~\ref{def:coupled-orthogonal}), that operand's tile
    does not change across consecutive iterations and is thus fully
    reused, that operand is \emph{stationary}.
    In the standard 2D convolution, where each dimension is orthogonal
    to exactly one of the three operands (Table~\ref{tab:coupling-conv}),
    one can make at most one operand stationary per memory level.
    The three classical dataflows arise from this choice:
    \begin{itemize}
      \item \emph{Weight-stationary}: innermost loops iterate over
        dimensions orthogonal to the weights (e.g., $P$, $Q$, $N$),
        keeping the weight tile resident.
      \item \emph{Input-stationary}: innermost loops iterate over
        dimensions orthogonal to the inputs (e.g., $M$), keeping the
        input tile resident.
      \item \emph{Output-stationary}: innermost loops iterate over
        dimensions orthogonal to the outputs (e.g., $C$, $R$, $S$),
        keeping the partial-sum tile resident and accumulating into it.
    \end{itemize}

  \item \textbf{Halo reuse (temporal, affine).}
    When the innermost iterated dimension is \emph{affine} for an
    operand consecutive tiles
    \emph{overlap} due to the sliding-window access pattern.
    The size of the overlap depends on the affine coefficients (strides
    and dilations) and the tile sizes.

    Halo reuse is only possible when none of the dimensions involved in
    the sum of indices are spatially parallelized at any level below
    the current one; otherwise, the overlapping data would reside in a
    different hardware instance and cannot be reused.

  \item \textbf{Spatial reuse.}
    On a \emph{fanout} level, loop order is irrelevant, only the
    \emph{choice of parallelized dimensions} matters.
    Parallelizing a dimension orthogonal to an operand enables
    \emph{multicast}; parallelizing a dimension coupled to the output
    enables \emph{reduction}.
    With affine dimensions, overlapping tiles across instances create
    additional multicast or reduction opportunities when selective
    support is available.
\end{enumerate}


% ------------------------------------------------------------
\subsection{Interaction of the Three Degrees of Freedom}
\label{sec:dof-interaction}
% ------------------------------------------------------------

Tiling, parallelism, and loop ordering are not independent in their
effect on performance:

\begin{itemize}
  \item \textbf{Tiling + parallelism} jointly determine what data must
    be where and when.
    The tile sizes at each level define the storage footprint, while the
    parallelism strategy determines how many copies of inner levels
    concurrently share an outer memory's bandwidth.

  \item \textbf{Loop ordering} then schedules the temporal sequence of
    data transfers and decides which reuse mechanism fires at each
    memory level.
    A specific loop order can only exploit stationarity or halo reuse
    if the tiling has allocated enough iterations to the relevant
    dimension at that level.
\end{itemize}

At any memory level, the permutation of the loops determines which
reuse fires:
\begin{enumerate}
  \item \emph{Stationarity} is locked by the set of innermost loops
    that are orthogonal to an operand.
  \item \emph{Halo reuse} depends on the position of the first
    \emph{coupled} loop and whether it involves an affine dimension.
  \item On \emph{spatial} levels, loop order does not change behavior;
    only the choice of parallelized dimensions matters.
\end{enumerate}

\noindent
This interplay means that evaluating a mapping requires considering all
three decisions simultaneously, which is precisely why the mapping
problem is so challenging:
the search space is the Cartesian product of all tiling configurations,
all parallelism assignments, and all loop permutations at every level.
Section~\ref{sec:bg-mapspace} quantifies the resulting combinatorial
explosion.
